% !TEX program = pdflatex
\documentclass[11pt,a4paper]{article}

% =========================================================================
% POLÍTICAS DEL PROYECTO Y GOBERNANZA
% Documento común a los proyectos CasaScore y Horus Sentinel
% Gestión de Proyectos TICS (EST315) -- Semestre 2026-II -- Team Sobrevivientes
% =========================================================================

% --- PAQUETES DE CONFIGURACIÓN ---
\usepackage[utf8]{inputenc}
\usepackage[T1]{fontenc}
\usepackage[spanish,es-tabla]{babel}
% Compatibilidad con babel-spanish reciente (flechas de TikZ y signo de porcentaje)
\AtBeginDocument{\shorthandoff{<>}\spanishplainpercent}
\usepackage{geometry}
\geometry{
    top=3.6cm,
    bottom=2.3cm,
    left=2.2cm,
    right=2.2cm,
    % PARTE MODIFICABLE: headheight debe ser mayor que la altura del logo más alto (2.56cm = 73pt).
    headheight=80pt,
    headsep=15pt,
    footskip=20pt
}

\usepackage{graphicx}
\usepackage{fancyhdr}
\usepackage[table]{xcolor}
\usepackage{titlesec}
\usepackage{booktabs}
\usepackage{tabularx}
\usepackage{array}
\usepackage{amsmath}
\usepackage{enumitem}
\usepackage{microtype}
\usepackage{hyperref}

% Logos: funciona a 2 o 3 niveles de profundidad dentro de docs/latex/
\graphicspath{{../../assets/img/}{../../../assets/img/}}

% --- PALETA INSTITUCIONAL FIEI / UNA PUNO ---
\definecolor{AzulUNA}{RGB}{10, 45, 90}
\definecolor{DoradoFIEI}{RGB}{180, 130, 20}
\definecolor{GrisOscuro}{RGB}{50, 50, 50}
\definecolor{AzulClaro}{RGB}{228, 236, 246}

% --- CONFIGURACIÓN DE LISTAS ---
\setlist{topsep=2pt, itemsep=2pt, parsep=1pt}

% --- COLUMNAS DE TABLA ---
\newcolumntype{L}{>{\raggedright\arraybackslash}X}
\newcolumntype{P}[1]{>{\raggedright\arraybackslash}p{#1}}
\newcolumntype{C}[1]{>{\centering\arraybackslash}p{#1}}
\renewcommand{\arraystretch}{1.25}

% --- FORMATO DE ENLACES ---
\hypersetup{
    colorlinks=true,
    linkcolor=AzulUNA,
    urlcolor=AzulUNA,
    citecolor=AzulUNA,
    pdftitle={Políticas del Proyecto y Gobernanza},
    pdfauthor={Team Sobrevivientes}
}

% --- FORMATO DE TÍTULOS ---
\titleformat{\section}{\color{AzulUNA}\normalfont\large\bfseries}{\thesection.}{0.7em}{}[\color{AzulUNA}\titlerule]
\titleformat{\subsection}{\color{AzulUNA}\normalfont\normalsize\bfseries}{\thesubsection.}{0.5em}{}
\titlespacing*{\section}{0pt}{9pt}{4pt}
\titlespacing*{\subsection}{0pt}{6pt}{2pt}

% --- ENCABEZADO Y PIE DE PÁGINA (fancyhdr) ---
\pagestyle{fancy}
\fancyhf{}

% ===== TAMAÑO DE LOS LOGOS (PARTE MODIFICABLE) =====
% logo_fiei.png  : height=1.45cm (lado izquierdo)
% logo_equipo.png: height=2.56cm (lado derecho; 25 % más grande porque la imagen tiene más borde blanco)
% Para cambiar el tamaño, modifique solo el valor height=... de cada \includegraphics.
% Si el logo más alto supera 2.6cm, aumente también headheight y top en la configuración de página.
\fancyhead[L]{%
    \includegraphics[height=1.45cm,keepaspectratio]{logo_fiei.png}%
}
\fancyhead[C]{%
    \scriptsize \textbf{\textcolor{AzulUNA}{UNIVERSIDAD NACIONAL DEL ALTIPLANO DE PUNO}}\\[2pt]
    \tiny \textcolor{GrisOscuro}{FACULTAD DE INGENIERÍA ESTADÍSTICA E INFORMÁTICA}\\[2pt]
    \scriptsize \textbf{\textcolor{AzulUNA}{Gestión de Proyectos TICS (EST315) \textbar\ Team Sobrevivientes}}\\[2pt]
    \footnotesize \textbf{\textcolor{DoradoFIEI}{POLÍTICAS DEL PROYECTO Y GOBERNANZA}}%
}
\fancyhead[R]{%
    \includegraphics[height=2.56cm,keepaspectratio]{logo_equipo.png}%
}

\renewcommand{\headrulewidth}{0pt}
\renewcommand{\headrule}{}

\fancyfoot[L]{\scriptsize \textcolor{GrisOscuro}{CasaScore \textperiodcentered\ Horus Sentinel}}
\fancyfoot[C]{\small \thepage}
\fancyfoot[R]{\scriptsize \textcolor{GrisOscuro}{Puno, Perú -- Semestre 2026-II}}
\renewcommand{\footrulewidth}{0.4pt}

\begin{document}

% =========================================================================
% CABECERA DEL DOCUMENTO
% =========================================================================

\begin{center}
    {\Large \textbf{\textcolor{AzulUNA}{POLÍTICAS DEL PROYECTO Y GOBERNANZA}}}\\[3pt]
    {\normalsize \textbf{Normas de trabajo del equipo para los proyectos CasaScore y Horus Sentinel}}\\[4pt]
    \footnotesize \textbf{Fecha de aprobación:} viernes 04/09/2026 (semana 2) \quad \textbar\quad \textbf{Versión:} 1.0
\end{center}

\vspace{0.1cm}

\noindent\begin{tabularx}{\textwidth}{@{}P{3.8cm} L@{}}
\toprule
\textbf{Curso:} & Gestión de Proyectos TICS (EST315) -- Semestre 2026-II \\
\textbf{Docente:} & Ing. Renzo Apaza Cutipa \\
\textbf{Equipo:} & Team Sobrevivientes \\
\textbf{Proyectos:} & \textbf{CasaScore (AlquilaPuno)}: recomendación multicriterio de cuartos en alquiler para estudiantes de la UNA Puno. \newline
\textbf{Horus Sentinel}: mapa predictivo de zonas de riesgo delictivo y sugerencia de rutas seguras. \\
\textbf{Alcance del documento:} & Normas comunes a ambos proyectos y políticas de producto de cada uno. \\
\bottomrule
\end{tabularx}

% =========================================================================
\section{Fundamento Metodológico}
% =========================================================================
\begin{itemize}[leftmargin=*]
    \item \textbf{Principio rector:} \textit{sistematizar antes de automatizar}. Ninguna tarea técnica se inicia sin insumos validados, un responsable asignado y un criterio de aceptación medible (ciclo Insumos $\rightarrow$ Proceso $\rightarrow$ Entregable).
    \item \textbf{Guía metodológica:} PM4DEV. Cada proyecto recorre las seis fases del ciclo (Inicio, Planificación, Implementación, Monitoreo, Adaptación y Cierre) y aplica los nueve procesos de gestión: alcance, cronograma, presupuesto y calidad como restricciones; equipo, interesados, información, riesgos y contratos como soporte.
    \item \textbf{Mejora continua:} cada proceso sigue el ciclo Planificar -- Hacer -- Verificar -- Adaptar.
    \item \textbf{Enfoque de desarrollo:} híbrido. Gestión predictiva para el alcance, el cronograma, el presupuesto y los riesgos; Scrum con sprints semanales para construir el software.
    \item \textbf{Marco de referencia del curso:} principios y dominios de desempeño del PMBOK.
\end{itemize}

% =========================================================================
\section{Gestión de la Cartera de Proyectos}
% =========================================================================

\noindent\begin{tabularx}{\textwidth}{@{}P{3.4cm} L L@{}}
\toprule
\rowcolor{AzulClaro}
\textbf{Aspecto} & \textbf{CasaScore} & \textbf{Horus Sentinel} \\
\midrule
Producto al cierre de la Unidad I & MVP móvil Android funcional. & Prototipo funcional con datos simulados. \\
Producto al cierre del semestre & MVP con registro e inicio de sesión, datos en la nube y mapa de cuartos. Horizonte del proyecto: 3 años. & MVP funcional: mapa de calor de riesgo y sugerencia de ruta segura. \\
Documentación & Expediente propio (entregables 01 a 28). & Expediente propio (entregables 01 a 28). \\
Dedicación del equipo & 8 h por semana por integrante. & 8 h por semana por integrante. \\
\bottomrule
\end{tabularx}

\begin{itemize}[leftmargin=*]
    \item Cada proyecto tiene su propia línea base de alcance, cronograma y presupuesto, y su propio registro de riesgos.
    \item Los documentos de gobernanza del equipo son comunes a ambos proyectos.
    \item Ningún recurso, entregable o cambio de un proyecto se carga al otro sin una solicitud de cambio aprobada.
\end{itemize}

% =========================================================================
\section{Organización del Equipo y Roles}
% =========================================================================
Los roles son los mismos en ambos proyectos.

\vspace{0.1cm}
{\small
\noindent\begin{tabularx}{\textwidth}{@{}P{3.9cm} P{3.4cm} L@{}}
\toprule
\rowcolor{AzulClaro}
\textbf{Integrante} & \textbf{Rol} & \textbf{Responsabilidades} \\
\midrule
\textbf{Paricahua Pari, Clyde Neil} & Project Manager & Dirección de ambos proyectos, control del alcance, aprobación de entregables y cambios menores, relación con el docente, decisión final ante desacuerdos y firma de actas. \\
\addlinespace[3pt]
\textbf{Apaza Solis, Ruth Karina} & Analista de Planificación & EDT, cronograma, presupuesto, tablero Kanban, cálculo del valor ganado (EVM) y control de la línea base. \\
\addlinespace[3pt]
\textbf{Tarqui Ajahuana, Eliseo} & Especialista Técnico / Arquitectura & Arquitectura, modelo de datos, algoritmos, desarrollo del software, pruebas técnicas y contratos con servicios tecnológicos. \\
\addlinespace[3pt]
\textbf{Quisbert Quispe, Viviana Mary} & Gestora de Comunicación, Riesgos y Documentación & Registro de riesgos, plan de información y comunicaciones, actas de reunión, custodia del repositorio documental y compilación en \LaTeX. \\
\bottomrule
\end{tabularx}
}

\subsection{Toma de decisiones}
\begin{itemize}[leftmargin=*]
    \item Las decisiones técnicas las propone el Especialista Técnico y las aprueba el Project Manager.
    \item Los cambios a una línea base siguen el proceso de solicitud de cambio del Plan de Dirección de cada proyecto.
    \item Los desacuerdos se resuelven primero por deliberación en reunión. Si persisten, decide el Project Manager y, cuando afectan el alcance aprobado, en coordinación con el docente.
    \item Cada integrante responde por su rol, pero todos colaboran en las tareas del sprint cuando el avance lo requiere.
\end{itemize}

% =========================================================================
\section{Protocolo de Reuniones}
% =========================================================================
\begin{itemize}[leftmargin=*]
    \item \textbf{Reuniones ordinarias:} martes y viernes de 20:00 a 21:00. El martes se controla el avance de ambos proyectos; el viernes se hace la revisión del sprint.
    \item \textbf{Agenda fija:} acuerdos de la reunión anterior, avance frente al tablero, impedimentos, riesgos y asignación de tareas.
    \item \textbf{Actas:} cada reunión se registra en un acta con avance real, desviaciones, acuerdos y compromisos. La redacta la Gestora de Comunicación y se publica en el repositorio.
    \item \textbf{Sesiones extraordinarias:} se convocan ante retrasos que afecten la ruta crítica o antes de entregas importantes. Las jornadas nocturnas se limitan a las semanas previas a una entrega, se registran en acta e incluyen retorno seguro de los integrantes.
    \item \textbf{Asistencia:} la inasistencia se comunica antes de la reunión y el integrante revisa el acta para asumir sus compromisos.
\end{itemize}

% =========================================================================
\section{Gestión de la Información y del Repositorio}
% =========================================================================
\begin{itemize}[leftmargin=*]
    \item \textbf{Repositorio único:} todo el código y la documentación de ambos proyectos se guarda en un solo repositorio de GitHub. Se prohíbe intercambiar archivos por canales no auditables.
    \item \textbf{Carpetas:} \texttt{docs/latex/00\_comun/} para los documentos del equipo, \texttt{docs/latex/casascore/} y \texttt{docs/latex/horus\_sentinel/} para cada proyecto, y \texttt{src/} para el código.
    \item \textbf{Numeración:} el número del entregable aparece solo en el nombre de la carpeta y del archivo. Dentro del documento, los demás documentos se citan por su nombre.
    \item \textbf{Formato:} todo documento formal se redacta en \LaTeX\ con el encabezado institucional y se acompaña de su PDF compilado. Los archivos Markdown se usan solo como borradores.
    \item \textbf{Fechas:} se indican con la semana del semestre (la semana 1 no se cuenta como semana de trabajo).
\end{itemize}

% =========================================================================
\section{Control de Versiones en GitHub}
% =========================================================================

\subsection{Ramas}
\begin{enumerate}[leftmargin=*]
    \item \textbf{Rama \texttt{main} protegida:} no se permiten envíos directos.
    \item \textbf{Ramas de trabajo con el nombre del proyecto:}
    \begin{itemize}
        \item \texttt{feature/<proyecto>-<tarea>}: nuevas funciones.
        \item \texttt{fix/<proyecto>-<error>}: corrección de defectos.
        \item \texttt{docs/<proyecto>-<documento>}: documentos (\texttt{comun} para los documentos del equipo).
    \end{itemize}
    \item \textbf{Pull Request obligatorio:} toda integración a \texttt{main} requiere la aprobación de al menos otro integrante.
\end{enumerate}

\subsection{Mensajes de confirmación (commits)}
Formato \texttt{tipo(proyecto-módulo): descripción}. Ejemplos:
\begin{itemize}[leftmargin=*]
    \item \texttt{feat(casascore-ranking): normalizacion min-max del precio}
    \item \texttt{feat(horus-mapa): mapa de calor con datos simulados}
    \item \texttt{fix(casascore-auth): validacion de credenciales}
    \item \texttt{docs(comun-politicas): version 1.0 de las politicas}
\end{itemize}

\subsection{Seguimiento}
\begin{itemize}[leftmargin=*]
    \item \textbf{Tablero Kanban} en GitHub Projects con los estados \textit{Por hacer}, \textit{En progreso}, \textit{En revisión} y \textit{Hecho}, y una etiqueta por proyecto (\texttt{casascore}, \texttt{horus}, \texttt{comun}). Lo administra la Analista de Planificación.
    \item \textbf{Issues:} cada entregable y paquete de trabajo tiene un Issue con un solo responsable, fecha límite y criterio de aceptación.
    \item \textbf{Versiones:} etiquetas por proyecto, por ejemplo \texttt{casascore-v1.0-mvp} y \texttt{horus-v0.1-prototipo}.
\end{itemize}

% =========================================================================
\section{Políticas de Calidad}
% =========================================================================
\begin{itemize}[leftmargin=*]
    \item Todo entregable tiene un criterio de aceptación definido antes de empezar a elaborarlo.
    \item Todo documento y todo código pasa por la revisión de otro integrante antes de integrarse.
    \item Los defectos se registran como Issues y se corrigen antes de cerrar el sprint.
    \item La calidad la juzga el usuario final: cada producto se valida con pruebas de usuarios antes de su aceptación.
\end{itemize}

% =========================================================================
\section{Protección de Datos y Confidencialidad}
% =========================================================================
\begin{itemize}[leftmargin=*]
    \item Los datos personales se tratan conforme a la Ley N.\textsuperscript{o} 29733, Ley de Protección de Datos Personales: con consentimiento, solo para la finalidad del proyecto y con medidas de seguridad.
    \item Se prohíbe vender, ceder o publicar datos de usuarios.
    \item No se guardan contraseñas, claves de acceso ni datos personales reales dentro del repositorio.
\end{itemize}

% =========================================================================
\section{Políticas de Producto}
% =========================================================================

\subsection{CasaScore (AlquilaPuno)}
\begin{itemize}[leftmargin=*]
    \item \textbf{MVP primero:} se priorizan usuarios, publicaciones, búsqueda, filtros y el motor del puntaje. El aprendizaje automático queda fuera del alcance.
    \item \textbf{Normalización 0--100:} ninguna variable (precio, distancia, seguridad, servicios) entra al motor sin llevarse antes a la escala común de 0 a 100.
    \item \textbf{Transparencia:} todo puntaje muestra los criterios que lo explican; no hay recomendaciones de caja negra ni posiciones pagadas.
    \item \textbf{Preferencias del usuario:} los pesos que define el estudiante prevalecen sobre cualquier valor por defecto.
    \item \textbf{Publicación completa:} un cuarto solo se publica con 4 fotos reales, precio, área, distancia, tiempo a pie y condiciones de convivencia declaradas.
    \item \textbf{Límites:} el sistema no cobra ni intermedia pagos, no firma contratos y no reemplaza la visita al inmueble.
    \item \textbf{Accesibilidad:} formularios simples, pensados para el celular y para usuarios sin conocimientos técnicos.
\end{itemize}

\subsection{Horus Sentinel}
\begin{itemize}[leftmargin=*]
    \item \textbf{Datos por etapas:} en la Unidad I el prototipo usa solo datos simulados. Para el MVP final se usan datos abiertos o datos autorizados por escrito por la institución que los entrega.
    \item \textbf{Datos agregados:} el sistema trabaja con incidentes agregados por zona y horario; no identifica personas, no usa cámaras ni reconocimiento facial.
    \item \textbf{Prevención del sesgo:} el modelo se audita para no sobrerrepresentar zonas históricamente más vigiladas, y se combinan fuentes oficiales con reportes ciudadanos.
    \item \textbf{Transparencia:} el mapa indica la fuente y la fecha de los datos, y el nivel de riesgo se presenta como una estimación, no como una certeza.
    \item \textbf{Alcance piloto:} se trabaja en un solo distrito piloto de Puno; no hay cobertura nacional.
    \item \textbf{Uso:} la herramienta apoya la prevención y la elección de rutas; no reemplaza la labor de la Policía ni del serenazgo.
\end{itemize}

% =========================================================================
\section{Suscripción y Conformidad}
% =========================================================================
\noindent En señal de conformidad con las presentes políticas, los integrantes del equipo firman a continuación:

\vspace{1.6cm}

\begin{center}
\begin{tabular}{c c c}
    \parbox{6.5cm}{\centering \rule{6.0cm}{0.6pt}\\[5pt] \textbf{Clyde Neil Paricahua Pari}\\[2pt] {\footnotesize Project Manager}}
    & \hspace{0.8cm} &
    \parbox{6.5cm}{\centering \rule{6.0cm}{0.6pt}\\[5pt] \textbf{Ruth Karina Apaza Solis}\\[2pt] {\footnotesize Analista de Planificación}}
    \\[2.0cm]
    \parbox{6.5cm}{\centering \rule{6.0cm}{0.6pt}\\[5pt] \textbf{Eliseo Tarqui Ajahuana}\\[2pt] {\footnotesize Especialista Técnico / Arquitectura}}
    & \hspace{0.8cm} &
    \parbox{6.5cm}{\centering \rule{6.0cm}{0.6pt}\\[5pt] \textbf{Viviana Mary Quisbert Quispe}\\[2pt] {\footnotesize Gestora de Comunicación, Riesgos y Documentación}}
\end{tabular}
\end{center}

\end{document}
